% Part: first-order-logic
% Chapter: natural-deduction
% Section: propositional-rules

\documentclass[../../../include/open-logic-section]{subfiles}

\begin{document}

\iftag{FOL}
      {\olfileid{fol}{ntd}{prl}}
      {\olfileid{pl}{ntd}{prl}}

\olsection{বচনমূলক বিধি}


\subsection{$\land$-এর বিধি}

\begin{defish}
\AxiomC{$!A$}
\AxiomC{$!B$}
\RightLabel{\Intro{\land}}
\BinaryInfC{$!A \land !B$}
\DisplayProof
\hfill
\begin{tabular}{r}
\AxiomC{$!A \land !B$}
\RightLabel{\Elim{\land}}
\UnaryInfC{$!A$}
\DisplayProof
\\[3ex]
\AxiomC{$!A \land !B$}
\RightLabel{\Elim{\land}}
\UnaryInfC{$!B$}
\DisplayProof
\end{tabular}
\end{defish}

\subsection{$\lor$-এর বিধি}

\begin{defish}
\begin{tabular}{r}
\AxiomC{$!A$}
\RightLabel{\Intro{\lor}}
\UnaryInfC{$!A \lor !B$}
\DisplayProof
\\[3ex]
\AxiomC{$!B$}
\RightLabel{\Intro{\lor}}
\UnaryInfC{$!A \lor !B$}
\DisplayProof
\end{tabular}
\hfill
\AxiomC{$!A \lor !B$}
\AxiomC{$\Discharge{!A}{n}$}
\DeduceC{$!C$}
\AxiomC{$\Discharge{!B}{n}$}
\DeduceC{$!C$}
\DischargeRule{\Elim{\lor}}{n}
\TrinaryInfC{$!C$}
\DisplayProof
\end{defish}

\subsection{$\lif$-এর বিধি}

\begin{defish}
\AxiomC{$\Discharge{!A}{n}$}
\DeduceC{$!B$}
\DischargeRule{\Intro{\lif}}{n}
\UnaryInfC{$!A \lif !B$}
\DisplayProof
\hfill
\AxiomC{$!A \lif !B$}
\AxiomC{$!A$}
\RightLabel{\Elim{\lif}}
\BinaryInfC{$!B$}
\DisplayProof
\end{defish}

\subsection{$\lnot$-এর বিধি}

\begin{defish}
\AxiomC{$\Discharge{!A}{n}$}
\noLine
\DeduceC{$\lfalse$}
\DischargeRule{\Intro{\lnot}}{n}
\UnaryInfC{$\lnot !A$}
\DisplayProof
\hfill
\AxiomC{$\lnot !A$}
\AxiomC{$!A$}
\RightLabel{\Elim{\lnot}}
\BinaryInfC{$\lfalse$}
\DisplayProof
\end{defish}

\subsection{$\lfalse$-এর বিধি}

\begin{defish}
\AxiomC{$\lfalse$}
\RightLabel{\FalseInt}
\UnaryInfC{$!A$}
\DisplayProof
\hfill
\AxiomC{$\Discharge{\lnot !A}{n}$}
\DeduceC{$\lfalse$}
\DischargeRule{\FalseCl}{n}
\UnaryInfC{$!A$}
\DisplayProof
\end{defish}

লক্ষ করো, $\Intro{\lnot}$ ও $\FalseCl$ খুবই সদৃশ। পার্থক্য হলো,
$\Intro{\lnot}$ একটি নেতিবাচক !!{sentence}~$\lnot !A$ নিষ্পাদন
করে, কিন্তু $\FalseCl$ একটি সদর্থক !!{sentence}~$!A$ নিষ্পাদন করে।

যখনই কোনো বিধিতে কোনো অনুমিতি অবমুক্ত করা যেতে পারে বলে দেখানো হয়,
তখন সেটিকে অনুমতি হিসেবে ধরি, বাধ্যবাধকতা হিসেবে নয়। যেমন,
$\Intro{\lif}$ বিধিতে পূর্বধারণা~$!B$-এর !!{derivation}-এ~$!A$
রূপের শূন্য-সহ যেকোনো সংখ্যক অনুমিতি অবমুক্ত করতে পারি।

\end{document}
